% Work-MAT | 2. rocnik | pripravovana hodina 29
% Návaznost: uživatelské připravy.zip (hodiny 01–20), autorská učebnice na GitHubu.
% Realisticky.cz: didaktická inspirace v posloupnosti a typech úloh; zadání jsou vlastní.
\documentclass[11pt,a4paper]{article}
\input{../../sablony/pripravy/priprava.sty}
\nastavPripravu{2. ročník}{Ověřování znalostí}{29}{Čtvrtletní písemná práce P1}
\begin{document}
\blok{Motivace a návaznost}
Čtvrtletní písemná práce uzavírá funkce a exponenciální rovnice před začátkem logaritmů. Ověř aktuální pravidla klasifikace v Moodle a schválené zadání.
\poznamka{Orientační nabídka pro 45--55 minut: krátké zopakování a motivace, nový výklad, tři společné vzorové úlohy, samostatná práce. Vyber podle tempa třídy; zbývající zadání slouží jako rezerva.}
\blok{Výklad}
Nejde o hotový test. Rozborové příklady níže jsou analogické úlohy pro pozdější zpětnou vazbu; bez schválení je nepoužívej jako skutečnou písemnou práci.
\blok{Společné příklady}
\begin{prikladReseniBox}[title={Příklad 1}]
Vyřeš, resp. urči: $3^x+3^{x+1}=108$.
\tcblower
\[\begin{aligned}4\cdot3^x=108 \\ 3^x=27\;\Longrightarrow\;x=3\end{aligned}\]
\end{prikladReseniBox}
\begin{prikladReseniBox}[title={Příklad 2}]
Vyřeš, resp. urči: $9^x-4\cdot3^x+3=0$.
\tcblower
\[\begin{aligned}t=3^x>0,\quad (t-1)(t-3)=0 \\ K=\{0,1\}\end{aligned}\]
\end{prikladReseniBox}
\poznamka{Po práci zaznamenej typické chyby; zadání testu ani bodování nevytvářej automaticky.}
\blok{Samostatné procvičování}
\begin{zadaniReseniBox}[title={Samostatná práce 1}]
\begin{enumerate}[label=\alph*)]
\item $4^x-5\cdot2^x+4=0$.
\item $2^{x+1}+2^x=24$.
\item $2^{2x-1}>2^{x+2}$.
\item $2^x+2^{3-x}=6$.
\item Proč $a^x$ nemůže být nulové?.
\end{enumerate}
\tcblower
\begin{enumerate}[label=\alph*)]
\item $K=\{0,2\}$.
\item $x=3$.
\item $x>3$.
\item $K=\{1,2\}$.
\item Obor hodnot $(0;\infty)$..
\end{enumerate}
\end{zadaniReseniBox}
\blok{Dobrovolné rozšíření}
\blok{Úkol na doma}
Při početních úlohách napiš postup, zvolenou metodu a zkoušku v původním zadání. U soustav kontroluj obě rovnice; u nerovnic výsledek zakresli na osu.
\begin{zadaniReseniBox}[title={Domácí úkol}]
\begin{enumerate}[label=\alph*)]
\item $3^{x+1}+3^x=108$.
\item $9^x+2\cdot3^x-15=0$.
\item $2^{x+2}-2^x=12$.
\item $2^{2x-1}>2^{x+2}$.
\item $2^x+2^{3-x}=6$.
\item Sepiš dvě hlavní početní metody..
\end{enumerate}
\tcblower
\begin{enumerate}[label=\alph*)]
\item $x=3$.
\item $x=1$.
\item $x=2$.
\item $x>3$.
\item $K=\{1,2\}$.
\item Vytýkání a substituce..
\end{enumerate}
\end{zadaniReseniBox}
\poznamka{Domácí úkol obsahuje 6 položek; předpokládaná domácí práce přibližně 15--30 minut. Po větších blocích jsou vědomě promíchané dříve probrané typy.}
\end{document}
